\documentclass[11pt]{rapportECL}
\usepackage{lipsum}
\usepackage{titlesec}
\titleformat{\section}{\normalfont\normalsize\bfseries}{\thesection}{1em}{}
\titleformat{\subsection}{\normalfont\normalsize\bfseries}{\thesubsection}{1em}{}
\titleformat{\subsubsection}{\normalfont\normalsize\bfseries}{\thesubsubsection}{1em}{}
\titlespacing*{\section}{0pt}{1ex plus .2ex}{.5ex plus .2ex}
\titlespacing*{\subsection}{0pt}{1ex plus .2ex}{.5ex plus .2ex}
\titlespacing*{\subsubsection}{0pt}{1ex plus .2ex}{.5ex plus .2ex}
\title{Note_de_lecture_TONGUE_KEVIN}
\begin{document}
\titre{Note de Lecture}
\UE{SHS-Management S9}
\sujet{Autonomie, apprentissage et responsabilité dans les organisations contemporaines}
\enseignant{Patrick \textsc{LAURENT}}
\eleves{Kevin \textsc{TONGUE}}
\UEU{ENISE-5GM}
% \fairemarges
\fairepagedegarde
\tabledematieres
\fontsize{11}{13}\selectfont

\section{Introduction}
Cette note de lecture combine les analyses de cinq textes relatifs à l'EC « SHS Management S9 ». Chaque section présente un résumé des idées principales suivies d'un commentaire argumenté, en lien avec les enseignements du cours.

\section{Note de Lecture : « L’autonomie ne se décrète pas, elle s’apprend ! Une approche des processus d’autonomisation des salariés à l’aune du concept de capabilité » de Benoît Grasser et Florent Noël (Revue Française de Gestion, n° 309, 2023)}

\subsection{Résumé des Idées et Positions Présentées dans le Texte}
Le texte mobilise le concept de capabilité de Sen pour analyser l’autonomie au travail, évitant les visions extrêmes taylorienne ou libérée. Via une étude de cas, il montre que l’autonomie accordée se traduit en initiatives si des facteurs de conversion (compétences, soutien) permettent des réalisations valorisées. C’est un processus dynamique avec tensions (émancipation vs. contrôle) et boucles d’apprentissage (activation, entretien, développement).

\subsection{Commentaire Argumenté}

\subsubsection{Définition et Justification d’une Problématique}
La problématique est : Dans quelle mesure l’autonomie au travail peut-elle être un processus d’apprentissage collectif intégrant tensions et facteurs de conversion pour des réalisations valorisées ? Elle est justifiée par la critique des approches binaires et s’appuie sur les théories des organisations (relations humaines, contingence).

\subsubsection{Définition et Justification d’un Plan d’Argumentation}
Le plan se structure en : (1) Processus dynamique ; (2) Rôle des tensions et facteurs de conversion ; (3) Boucles d’apprentissage.

\subsubsection{Déroulement Discursif du Plan d’Argumentation}
\begin{enumerate}
\item L’autonomie est un processus émergent, résonnant avec l’école classique critiquée et les théories de contingence (Blau) pour l’adaptation.
\item Les facteurs de conversion transforment ressources en capabilités, avec tensions ; aligné sur Williamson (coûts de transaction) et relations humaines (Mayo).
\item Les boucles d’apprentissage maintiennent l’autonomie ; s’appuie sur empiristes (Sloan) et contingence (Woodward).
\end{enumerate}

\subsubsection{Conclusion Persuasive}
L’autonomie émerge comme processus collectif, invitant à favoriser l’apprentissage pour une performance durable.

\section{Note de Lecture : « Leaders’ Critical Role in Building a Learning Culture » by Henrik Saabye and Thomas Borup (MIT Sloan Management Review, Spring 2025)}

\subsection{Résumé des Idées et Positions Présentées dans le Texte}
Le texte souligne le rôle des leaders comme facilitateurs d’apprentissage face au VUCA, critiquant la délégation externe. Via études chez Lego et Velux, il prône « aller lentement pour aller vite » avec méthodes comme A3 thinking pour développer la résolution de problèmes, créant un effet ripple.

\subsection{Commentaire Argumenté}

\subsubsection{Définition et Justification d’une Problématique}
La problématique est : Dans quelle mesure les leaders doivent-ils être facilitateurs d’apprentissage pour une culture adaptative ? Justifiée par les limites de la délégation et liée aux empiristes (Sloan) et contingence (Blau).

\subsubsection{Définition et Justification d’un Plan d’Argumentation}
Le plan : (1) Rôle contre-intuitif ; (2) Intégration des défis ; (3) Effet d’entraînement.

\subsubsection{Déroulement Discursif du Plan d’Argumentation}
\begin{enumerate}
\item Leaders priorisent l’apprentissage contextuel ; contredit école classique, aligne sur contingence (Blau).
\item Méthodes structurées transforment problèmes ; renforcé par empiristes (Sloan) et relations humaines (Bavelas).
\item Effet ripple crée culture durable ; s’appuie sur Williamson et contingence (Woodward).
\end{enumerate}

\subsubsection{Conclusion Persuasive}
Les leaders facilitent l’apprentissage pour résilience, appelant à un shift managérial aligné sur les théories du cours.

\section{Note de Lecture : « Management Innovation Made in China: Haier’s Rendanheyi » by Jedrzej George Frynas, Michael J. Mol, and Kamel Mellahi (California Management Review, Vol. 61, 2018)}

\subsection{Résumé des Idées et Positions Présentées dans le Texte}
Le texte explore Rendanheyi chez Haier pour gérer le VUCA, transformant hiérarchie en réseau entrepreneurial via micro-entreprises. Dépend de contextes (organisationnel, compétitif, etc.) ; propose modèle étendu pour innovations émergentes.

\subsection{Commentaire Argumenté}

\subsubsection{Définition et Justification d’une Problématique}
La problématique est : Dans quelle mesure Rendanheyi transforme-t-il les organisations via facteurs contextuels pour performance durable ? Liée à SCP (Bain) et contingence (Blau).

\subsubsection{Définition et Justification d’un Plan d’Argumentation}
Le plan : (1) Innovation contextuelle ; (2) Intégration des facteurs ; (3) Modèle étendu.

\subsubsection{Déroulement Discursif du Plan d’Argumentation}
\begin{enumerate}
\item Rendanheyi face au VUCA ; aligne sur contingence (Blau).
\item Facteurs transforment structure ; renforcé par Williamson et empiristes (Sloan).
\item Modèle lie structure-performance ; s’appuie sur SCP et relations humaines (Mayo).
\end{enumerate}

\subsubsection{Conclusion Persuasive}
Rendanheyi offre leçons pour adaptabilité, incitant à repenser pratiques.

\section{Note de Lecture : « Building Culture From the Middle Out » by Spencer Harrison and Kristie Rogers (MIT Sloan Management Review, Spring 2024)}

\subsection{Résumé des Idées et Positions Présentées dans le Texte}
Le texte met l’accent sur leaders intermédiaires pour enrichir culture (big-C/small-c) via endossement et enrichissement, basé sur étude Fortune 100. Crée adaptabilité avec bénéfices comme engagement.

\subsection{Commentaire Argumenté}

\subsubsection{Définition et Justification d’une Problématique}
La problématique est : Dans quelle mesure les intermédiaires construisent-ils culture en liant big-C/small-c pour adaptabilité ? Liée à empiristes (Sloan) et contingence (Blau).

\subsubsection{Définition et Justification d’un Plan d’Argumentation}
Le plan : (1) Distinction big-C/small-c ; (2) Stratégies intégrées ; (3) Mécanismes d’innovation.

\subsubsection{Déroulement Discursif du Plan d’Argumentation}
\begin{enumerate}
\item Cadre pour rôle actif ; aligne sur contingence (Blau).
\item Stratégies lient valeurs ; renforcé par empiristes (Schein) et relations humaines (Bavelas).
\item Innovation pour culture évolutive ; s’appuie sur Williamson et contingence (Woodward).
\end{enumerate}

\subsubsection{Conclusion Persuasive}
Intermédiaires enrichissent culture pour résilience, appelant à innovation.

\section{Note de Lecture : « Corporate Responsibility Meets the Digital Economy » by Leena Lankoski and N. Craig Smith (California Management Review, Vol. 67, 2025)}

\subsection{Résumé des Idées et Positions Présentées dans le Texte}
Le texte revisite CR face au digital via questions (quoi, envers qui, qui), élargissant enjeux, stakeholders et attribution. Implications pour managers, advocates, etc., pour équilibrer innovation-éthique.

\subsection{Commentaire Argumenté}

\subsubsection{Définition et Justification d’une Problématique}
La problématique est : Dans quelle mesure la digitalisation nécessite-t-elle revisite CR pour adaptation éthique ? Liée à économiques (Williamson) et empiristes (Schein).

\subsubsection{Définition et Justification d’un Plan d’Argumentation}
Le plan : (1) Revisite questions ; (2) Intégration impacts ; (3) Implications pratiques.

\subsubsection{Déroulement Discursif du Plan d’Argumentation}
\begin{enumerate}
\item Réponse au digital ; aligne sur contingence (Woodward).
\item Impacts équilibrent éthique ; renforcé par Williamson et empiristes (Schein).
\item Mécanismes pour adaptation ; s’appuie sur SCP et relations humaines (Bavelas).
\end{enumerate}

\subsubsection{Conclusion Persuasive}
CR digitale comme pilier, incitant à repenser pratiques pour résilience.

\section{Conclusion Générale}
Les textes interconnectent autonomie, apprentissage, innovation, culture et CR, soulignant approches adaptatives pour performance dans environnements complexes.

\end{document}